\begin{lemma}
\label{lemma-decompose-associated-complexes}
Let $\mathcal{A}$ be an abelian category.
Let $U$ be a simplicial object of $\mathcal{A}$.
The canonical morphism of chain complexes
$N(U) \to s(U)$ is split. In fact,
$$
s(U) = N(U) \oplus D(U)
$$
for some complex $D(U)$. The construction $U \mapsto D(U)$
is functorial.
\end{lemma}

\begin{proof}
Define $D(U)_n$ to be the image of
$$
\bigoplus\nolimits_{\varphi : [n] \to [m]\text{ surjective}, \ m < n} N(U_m)
\xrightarrow{\bigoplus U(\varphi)}
U_n
$$
which is a subobject of $U_n$ complementary to
$N(U_n)$ according to Lemma \ref{lemma-splitting-abelian-category} and
Definition \ref{definition-split}. We show that
$D(U)$ is a subcomplex. Pick a surjective
map $\varphi : [n] \to [m]$ with $m < n$ and consider
the composition
$$
N(U_m) \xrightarrow{U(\varphi)} U_n \xrightarrow{d_n} U_{n - 1}.
$$
This composition is the sum of the maps
$$
N(U_m) \xrightarrow{U(\varphi \circ \delta^n_i)} U_{n - 1}
$$
with sign $(-1)^i$, $i = 0, \ldots, n$.

\medskip\noindent
First we will prove by ascending induction on $m$,
$0 \leq m < n - 1$ that all the maps $U(\varphi \circ \delta^n_i)$
map $N(U_m)$ into $D(U)_{n - 1}$. (The case $m = n - 1$ is treated below.)
Whenever the map $\varphi \circ \delta^n_i : [n - 1] \to [m]$
is surjective then the image of $N(U_m)$ under $U(\varphi \circ \delta^n_i)$
is contained in $D(U)_{n - 1}$ by definition.
If $\varphi \circ \delta^n_i : [n - 1] \to [m]$ is not surjective,
set $j = \varphi(i)$ and observe that $i$ is the unique
index whose image under $\varphi$ is $j$. We may write
$\varphi \circ \delta^n_i = \delta^m_j \circ \psi \circ \delta^n_i$
for some $\psi : [n - 1] \to [m - 1]$. Hence
$U(\varphi \circ \delta^n_i) = U(\psi \circ \delta^n_i) \circ d^m_j$ which
is zero on $N(U_m)$ unless $j = m$. If $j = m$, then
$d^m_m(N(U_m)) \subset N(U_{m - 1})$ and hence
$U(\varphi \circ \delta^n_i)(N(U_m)) \subset
U(\psi \circ \delta^n_i)(N(U_{m - 1}))$ and we win
by induction hypothesis.

\medskip\noindent
To finish proving that $D(U)$ is a subcomplex
we still have to deal with the composition
$$
N(U_m) \xrightarrow{U(\varphi)} U_n \xrightarrow{d_n} U_{n - 1}.
$$
in case $m =  n - 1$. In this case $\varphi = \sigma^{n - 1}_j$
for some $0 \leq j \leq n - 1$ and $U(\varphi) = s^{n - 1}_j$.
Thus the composition is given by the sum
$$
\sum (-1)^i d^n_i \circ s^{n - 1}_j
$$
Recall from Remark \ref{remark-relations} that
$d^n_j \circ s^{n - 1}_j = d^n_{j + 1} \circ s^{n - 1}_j = \text{id}$
and these drop out because the corresponding terms have opposite signs.
The map $d^n_n \circ s^{n - 1}_j$, if $j < n - 1$, is equal to
$s^{n - 2}_j \circ d^{n - 1}_{n - 1}$. Since
$d^{n - 1}_{n - 1}$ maps $N(U_{n - 1})$ into $N(U_{n - 2})$,
we see that the image $d^n_n ( s^{n - 1}_j (N(U_{n - 1}))$
is contained in $s^{n - 2}_j(N(U_{n - 2}))$ which
is contained in $D(U_{n - 1})$ by definition. For all
other combinations of $(i, j)$ we have
either $d^n_i \circ s^{n - 1}_j = s^{n - 2}_{j - 1} \circ d^{n - 1}_i$
(if $i < j$), or
$d^n_i \circ s^{n - 1}_j = s^{n - 2}_j \circ d^{n - 1}_{i - 1}$
(if $n > i > j + 1$) and in these cases the map is zero because
of the definition of $N(U_{n - 1})$.
\end{proof}
